%% toc: P3 34 34.1,34.2,34.3,34.4
\chapter{Simple, NIP, and O-Minimal Theories}
\label{sec:p03-w03-p03-035-chapter}

\emph{Purpose.}  This chapter proves the combinatorial characterizations of
the independence property for one formula and for a complete theory, and the
implications among stability, simplicity, NIP and o-minimality that the
overview of \cref{sec:p03-nip-ominimal} records.  It then checks the standard
examples against the definitions: the random graph, dense linear orders and
real closed fields.

\emph{Inputs.}  A complete theory and the definitions of
\cref{def:p03-stability,def:p03-dividing-forking,def:p03-simple-theory,def:p03-nip,def:p03-o-minimal},
together with the ambient-model convention below.  The quantifier-elimination
results used for the examples are those of
\cref{thm:p03-rn-dlo-pattern-truth,thm:p03-w03-p03-015-graph-types} and
\cref{sec:p03-w03-p03-037-rcf}.

\emph{Conclusions.}  Equivalent forms of the independence property, with all
quantifiers explicit (\cref{thm:p03-w03-p03-035-ip-characterization}); the
dichotomy between exponential and polynomial counts of formula types
(\cref{thm:p03-w03-p03-035-sauer-shelah}); the reduction of theory-level NIP to
formulas in one variable (\cref{thm:p03-w03-p03-035-one-variable}); the
implications \emph{stable implies NIP}, \emph{unstable implies independence
property or strict order property}, \emph{simple and NIP implies stable},
\emph{o-minimal implies NIP}, \emph{stable implies local character of
dividing}, and \emph{stable implies simple}; and the position of the
examples, including an ordered theory that is NIP and not o-minimal.

\emph{Sources and versions.}  Shelah's 1971 Annals paper
\cite{Shelah1971Stability} is read in the published printing.  Simon's
\emph{Guide to NIP Theories} is read in the author-typeset version dated
16 October 2015 \cite{Simon2015GuideManuscript}, which the author distributes
with corrections to the 2015 printing \cite{Simon2015}; every locator to
Simon's text in this chapter refers to that manuscript.

\section{Formula-level independence and the NIP condition}
\label{sec:p03-w03-p03-035-nip}

\begin{hbconvention}{def:p03-w03-p03-035-ambient}{Ambient model, partitioned formulas and formula types}
Fix a complete theory \(T\) with infinite models and work in one
\(\bar\kappa\)-saturated model \(\mathfrak C\entails T\) as in
\cref{def:p03-w03-p03-031-monster}, with the additional requirement
\(\bar\kappa>2^{\abs T}\).  Thus sets of size at most \(2^{\abs T}\), in
particular sequences of length \(\abs T^+\) and countable sets of parameters,
are small, and every set of formulas in a small number of variables with
small parameter set that is finitely satisfied in \(\mathfrak C\) is realized
in \(\mathfrak C\) (\cref{thm:p03-w03-p03-031-saturation}).

A \emph{partitioned formula} \(\varphi(x;y)\) is an \(\Lang\)-formula without
parameters together with a division of its free variables into an object
tuple \(x\) and a parameter tuple \(y\), both finite.  For a tuple \(b\) of
length \(\abs y\), \(\varphi(\mathfrak C;b)\) is the set of \(\abs x\)-tuples
\(a\) with \(\mathfrak C\entails\varphi(a;b)\).  The \emph{opposite} formula
\(\varphi^{\mathrm{opp}}(y;x)\) is the same formula with the roles of \(x\)
and \(y\) exchanged.  Write \(\varphi^1=\varphi\) and
\(\varphi^0=\neg\varphi\).  For a set \(B\) of \(\abs y\)-tuples, a
\emph{\(\varphi\)-type over \(B\)} is a maximal set of formulas
\(\varphi(x;b)^{\varepsilon}\) with \(b\in B\) and
\(\varepsilon\in\{0,1\}\) that is realized in \(\mathfrak C\); their set is
\(S_\varphi(B)\).  An indexed family \((\varphi(x;b_i):i\in I)\) is
\emph{\(k\)-inconsistent} when for all pairwise distinct
\(i_1,\ldots,i_k\in I\) the conjunction
\(\bigwedge_{l=1}^{k}\varphi(x;b_{i_l})\) is realized by no tuple of
\(\mathfrak C\); repetitions among the \(b_i\) are counted by index.  These
conventions are those of
\cite[Section~1.1, pp.~10--11]{Simon2015GuideManuscript}, except that
Simon's monster model is also assumed homogeneous; no homogeneity is used
here, and every use of an automorphism in Simon's arguments is replaced by a
type-preserving copy.
\end{hbconvention}

A partitioned formula is a syntactic object paired with a choice of roles;
the same formula with another partition is a different partitioned formula,
which is why the opposite formula needs a name.  Formula types are needed
because type counting is done one formula at a time.

\begin{hbdefinition}{def:p03-w03-p03-035-shattering}{Shattering in the ambient model, VC-dimension and shatter counts}
Let \(\varphi(x;y)\) be a partitioned formula.  A set \(A\) of
\(\abs x\)-tuples of \(\mathfrak C\) is \emph{shattered by \(\varphi\)} when
for every \(A_0\subseteq A\) there is \(b\) with
\(\varphi(\mathfrak C;b)\cap A=A_0\).  The \emph{VC-dimension} of
\(\varphi\) is the largest \(n<\omega\) such that some set of \(n\) distinct
tuples is shattered by \(\varphi\), and is \(\infty\) when no largest such
\(n\) exists \cite[Remark~2.3, p.~14]{Simon2015GuideManuscript}.  For
\(n<\omega\), the \emph{shatter count} \(\pi_\varphi(n)\) is the maximum of
\(\abs{S_\varphi(B)}\) over sets \(B\) of exactly \(n\) distinct
\(\abs y\)-tuples.
\index{shattering}\index{VC-dimension}
\end{hbdefinition}

The definition of IP and NIP is that of \cref{def:p03-nip}: arbitrarily
large finite sets are shattered in some model of \(T\).  Simon instead
defines \(\varphi\) to be NIP when no infinite set is shattered
\cite[Definition~2.1, p.~13]{Simon2015GuideManuscript}.  The first theorem
shows that the two definitions agree and gives the other forms used below.

\begin{hbtheorem}{thm:p03-w03-p03-035-ip-characterization}{proved=annex:p03-w03-p03-035-ip-characterization}{Equivalent forms of the independence property}
For a partitioned formula \(\varphi(x;y)\) the following are equivalent.
\begin{enumerate}
  \item \(\varphi\) has IP in \(T\) in the sense of \cref{def:p03-nip}.
  \item For every \(n<\omega\) there is a set of \(n\) distinct
        \(\abs x\)-tuples of \(\mathfrak C\) shattered by \(\varphi\).
  \item There are pairwise distinct \(\abs x\)-tuples \((a_i:i<\omega)\) of
        \(\mathfrak C\) such that for all finite disjoint
        \(I,J\subseteq\omega\) some \(b\) satisfies \(\varphi(a_i;b)\) for
        every \(i\in I\) and \(\neg\varphi(a_j;b)\) for every \(j\in J\).
  \item \(\varphi^{\mathrm{opp}}(y;x)\) has IP in \(T\).
  \item There are an indiscernible sequence \((a_i:i<\omega)\) of
        \(\abs x\)-tuples and a tuple \(b\) such that, for every
        \(i<\omega\), \(\mathfrak C\entails\varphi(a_i;b)\) if and only if
        \(i\) is even.
  \item For every \(n<\omega\), \(\pi_\varphi(n)=2^n\).
\end{enumerate}
Moreover, in \textup{(3)} the set \(\{a_i:i<\omega\}\) is shattered by
\(\varphi\) in the sense of \cref{def:p03-w03-p03-035-shattering}.
\end{hbtheorem}

Item~(3) is Simon's ``infinite set shattered'' read through compactness;
item~(4) is \cite[Lemma~2.5, p.~14]{Simon2015GuideManuscript}; item~(5) is
\cite[Lemma~2.7, p.~15]{Simon2015GuideManuscript}, which Simon attributes to
Poizat \cite[p.~36]{Simon2015GuideManuscript}.  Item~(6) is the count of
formula types over finite sets; the next theorem says what happens when it
fails.

\begin{hbtheorem}{thm:p03-w03-p03-035-sauer-shelah}{proved=annex:p03-w03-p03-035-sauer-shelah}{Exponential or polynomial type counts}
Let \(\varphi(x;y)\) be a partitioned formula.
\begin{enumerate}
  \item \emph{(Sauer--Shelah bound.)}  If \(\varphi^{\mathrm{opp}}\) has
        finite VC-dimension \(d\), then for every \(n<\omega\)
        \[
          \pi_\varphi(n)\leq\sum_{i=0}^{d}\binom ni\leq(n+1)^d .
        \]
  \item \(\varphi\) is NIP if and only if there are \(d<\omega\) such that
        \(\pi_\varphi(n)\leq(n+1)^d\) for every \(n<\omega\), if and only if
        \(\pi_\varphi(n)<2^n\) for some \(n<\omega\).
\end{enumerate}
\end{hbtheorem}

The combinatorial bound is \cite[Lemma~6.4, p.~86]{Simon2015GuideManuscript},
stated there for an abstract family of sets; Simon's Example~6.2 on p.~86
identifies the VC-dimension of a formula with that of its family of
instances.  The proof in the annex is an induction on the size of the base
set different from Simon's compression argument.

\begin{hblemma}{thm:p03-w03-p03-035-em-extraction}{proved=annex:p03-w03-p03-035-em-extraction}{Indiscernible sequences realizing an Ehrenfeucht--Mostowski type}
Let \(A\) be a small set, \((a_i:i<\omega)\) a sequence of tuples of one
length, and \((X,<)\) a small infinite linear order.  There is an
\(A\)-indiscernible sequence \((c_t:t\in X)\) such that, for every
\(\Lang(A)\)-formula \(\theta(z_1,\ldots,z_n)\), if
\(\mathfrak C\entails\theta(a_{i_1},\ldots,a_{i_n})\) for all
\(i_1<\cdots<i_n<\omega\), then
\(\mathfrak C\entails\theta(c_{t_1},\ldots,c_{t_n})\) for all
\(t_1<\cdots<t_n\) in \(X\).
\end{hblemma}

The set of such formulas \(\theta\) is the Ehrenfeucht--Mostowski type of the
sequence over \(A\) \cite[Section~1.1.1, p.~11]{Simon2015GuideManuscript};
Simon refers the extraction to Ramsey's theorem and compactness, and the
annex writes that argument out.

\begin{hbtheorem}{thm:p03-w03-p03-035-alternation}{proved=annex:p03-w03-p03-035-alternation}{Bounded alternation on indiscernible sequences}
Let \(\varphi(x;y)\) be NIP.  There is \(N<\omega\) such that for every
indiscernible sequence \((a_i:i\in I)\) of \(\abs x\)-tuples and every
\(b\) there are no \(i_0<\cdots<i_N\) in \(I\) with
\(\mathfrak C\entails\varphi(a_{i_k};b)\leftrightarrow\neg\varphi(a_{i_{k+1}};b)\)
for all \(k<N\).  Consequently \(\varphi\) is NIP if and only if, for every
indiscernible sequence \((a_i:i\in I)\) with \(I\) nonempty and every
\(b\), there are a nonempty end segment \(I_0\subseteq I\) and
\(\varepsilon\in\{0,1\}\) with \(\mathfrak C\entails\varphi(a_i;b)^{\varepsilon}\)
for all \(i\in I_0\).
\end{hbtheorem}

This is \cite[Proposition~2.8, p.~16]{Simon2015GuideManuscript}; the least
such \(N\) bounds Simon's alternation number
\(\operatorname{alt}(\varphi)\) \cite[p.~15]{Simon2015GuideManuscript}.

\begin{hblemma}{thm:p03-w03-p03-035-boolean}{proved=annex:p03-w03-p03-035-boolean}{Boolean combinations of NIP formulas}
If \(\varphi(x;y)\) and \(\psi(x;y)\) are NIP, so are \(\neg\varphi\) and
\(\varphi\wedge\psi\).  Adding to \(y\) variables that do not occur in a
formula does not change whether it is NIP.  Hence every Boolean combination
of NIP formulas in the object variable \(x\) is NIP, its parameter variables
being the union of theirs.
\end{hblemma}

\begin{hbtheorem}{thm:p03-w03-p03-035-one-variable}{proved=annex:p03-w03-p03-035-one-variable}{NIP theories are detected by one-variable formulas}
The following are equivalent.
\begin{enumerate}
  \item \(T\) is NIP: every partitioned formula is NIP.
  \item Every partitioned formula \(\varphi(x;y)\) with \(\abs y=1\) is NIP.
  \item Every partitioned formula \(\varphi(x;y)\) with \(\abs x=1\) is NIP.
\end{enumerate}
If \(T\) is NIP, every formula \(\varphi(x;y,d)\) with parameters \(d\) from
\(\mathfrak C\) is NIP as a partitioned formula in \(x\) with parameter
variables \(y\).
\end{hbtheorem}

The equivalence of (1) and (2) is
\cite[Proposition~2.11, p.~16]{Simon2015GuideManuscript}; its proof there and
in the annex follows Poizat's argument through indiscernible sequences
\cite[p.~36]{Simon2015GuideManuscript}, not Shelah's original proof by
counting types, which Simon describes as using an absoluteness argument.
The last sentence is the remark after Simon's Definition~2.10 on p.~16.

\begin{hbwarning}{def:p03-w03-p03-035-formula-vs-theory}{One NIP formula against an NIP ambient theory}
\Cref{thm:p03-w03-p03-035-ip-characterization,thm:p03-w03-p03-035-sauer-shelah,thm:p03-w03-p03-035-alternation,thm:p03-w03-p03-035-boolean}
are statements about a single formula and hold in every complete theory,
including theories in which other formulas have IP.  Simon's text works at
this level only in Section~2.1; from there on the ambient theory is assumed
NIP, and Simon notes that many later results could be proved for one NIP
formula inside a theory with IP but are not
\cite[p.~9]{Simon2015GuideManuscript}.  So a result of Simon's later chapters
(invariant types, honest definitions, forking over models, measures) may
not be applied to one NIP formula of a theory with IP on the strength of the
text alone.  The implication results of
\cref{sec:p03-w03-p03-035-intersections} are statements about complete
theories and are labelled as such.
\end{hbwarning}

\begin{hbexample}{def:p03-w03-p03-035-simon-examples}{Simon's first examples of NIP and IP formulas}
These are the examples of \cite[Example~2.4, p.~14]{Simon2015GuideManuscript},
with the arguments.
\begin{itemize}
  \item In dense linear orders, language \(\{<\}\), the formula
        \(x\leq y\) has VC-dimension \(1\): a single point \(a\) is
        shattered by the parameters \(a\) and any point below \(a\); and if
        \(a_1<a_2\), no \(b\) satisfies \(a_2\leq b\) together with
        \(\neg(a_1\leq b)\), so \(\{a_2\}\) is not cut out of
        \(\{a_1,a_2\}\).
  \item In the complete theory of \((\mathbb N;+,\cdot)\), the formula
        ``\(x\) divides \(y\)'' has IP: for distinct primes
        \(p_0,\ldots,p_{N-1}\) and \(I\subseteq N\), the product
        \(\prod_{i\in I}p_i\) is divisible by \(p_i\) exactly when \(i\in I\).
  \item In the random graph, \(xEy\) has IP; every finite set of vertices is
        shattered (\cref{thm:p03-w03-p03-035-random-graph-ip}).
  \item In an infinite Boolean algebra in the language
        \(\{0,1,\wedge,\vee,\neg\}\), the formula \(x\leq y\), defined by
        \(x\wedge y=x\), has IP: a set of pairwise disjoint nonzero elements
        is shattered by the joins of its subsets, and an infinite Boolean
        algebra has arbitrarily large finite sets of pairwise disjoint
        nonzero elements (\cref{thm:p03-w03-p03-035-boolean-algebra-ip}).
\end{itemize}
Simon states the shattering of pairwise disjoint sets; the existence of
large pairwise disjoint sets in every infinite Boolean algebra, which the
example also needs, is proved in the annex to the next lemma.  The VC-dimension of a
formula and of its opposite can differ: in the theory of an infinite set
with equality only, \(x=y_1\vee x=y_2\vee x=y_3\) has VC-dimension \(3\)
and its opposite has VC-dimension \(2\)
\cite[Remark~2.6, pp.~14--15]{Simon2015GuideManuscript}; only finiteness is
preserved (\cref{thm:p03-w03-p03-035-ip-characterization}).
\end{hbexample}

\begin{hblemma}{thm:p03-w03-p03-035-boolean-algebra-ip}{proved=annex:p03-w03-p03-035-boolean-algebra-ip}{The order of an infinite Boolean algebra has IP}
Let \(B\) be an infinite Boolean algebra in the language
\(\{0,1,\wedge,\vee,\neg\}\), and write \(x\leq y\) for \(x\wedge y=x\).
\begin{enumerate}
  \item For every \(n<\omega\) there are pairwise distinct
        \(d_0,\ldots,d_{n-1}\in B\), all nonzero, with \(d_i\wedge d_j=0\)
        for \(i\neq j\).
  \item If \(d_0,\ldots,d_{n-1}\) are as in \textup{(1)} and, for
        \(I\subseteq n\), \(b_I=\bigvee_{i\in I}d_i\) (with \(b_\emptyset=0\)),
        then \(d_j\leq b_I\) holds exactly when \(j\in I\).
  \item Consequently \(\{d_0,\ldots,d_{n-1}\}\) is shattered in \(B\) by the
        partitioned formula \(x\leq y\), and \(x\leq y\) has IP in
        \(\operatorname{Th}(B)\) in the sense of \cref{def:p03-nip}.
\end{enumerate}
\end{hblemma}

The proof of \textup{(1)} splits an element with infinitely many elements
below it into two disjoint nonzero pieces, one of which again has infinitely
many elements below it, and collects the other pieces.  Every model of a
complete theory of an infinite Boolean algebra is an infinite Boolean
algebra, so the lemma applies to each of them; this is the setting of
Simon's item.

\subsection{Shelah's 1971 presentation of formula properties}
\label{sec:p03-w03-p03-035-shelah}

Shelah's \emph{Stability, the f.c.p., and superstability} is a long paper
whose \S\S 2--4 study properties of single formulas.  Its notation differs
from the one above.  Shelah writes \(\varphi(\bar x;\bar y)\), uses
\(S^m_\varphi(A)\) for complete \(\varphi\)-\(m\)-types over a set \(A\) of
elements (instances \(\varphi(\bar x;\bar a)\) with \(\bar a\) a tuple from
\(A\)), and writes \(\varphi^{\text{if }(\ldots)}\) for \(\varphi\) or
\(\neg\varphi\) according to whether the condition holds
\cite[\S 1 and pp.~296--297]{Shelah1971Stability}.

\begin{hbdefinition}{def:p03-w03-p03-035-shelah-defs}{Shelah's formula properties}
In \cite{Shelah1971Stability}:
\begin{enumerate}
  \item \(\varphi(\bar x;\bar y)\) is \emph{stable in \(\lambda\)} when
        \(\abs A\leq\lambda\) implies \(\abs{S^m_\varphi(A)}\leq\lambda\);
        it is \emph{stable} when stable in some cardinal; \(T\) is stable in
        \(\lambda\) when \(\abs A\leq\lambda\) implies
        \(\abs{S(A)}\leq\lambda\) (Definition~2.1, p.~297).
  \item \(\varphi(\bar x;\bar y)\) has the \emph{order property} when there
        are \(\bar a^0,\bar a^1,\ldots\) such that for every \(r<\omega\)
        the set \(\{\varphi(\bar x;\bar a^n)^{\text{if }(n\geq r)}:n<\omega\}\)
        is consistent (Definition~2.2, p.~297).
  \item \(\Gamma(\varphi,m,\alpha)\) is the set of formulas
        \(\varphi(\bar x_\eta;\bar y_{\eta\restriction k})^{\eta(k)}\) for
        \(\eta\in{}^{\alpha}2\) and \(k<\alpha\), together with \(T\); its
        consistency for \(\alpha=\omega\) is a binary tree of instances
        (Definition~2.3, p.~297).  The rank
        \(\operatorname{Rank}^m_\varphi\) of \(\varphi\)-\(m\)-types
        (Definitions~2.4--2.5, pp.~299--300) is
        \cref{def:p03-w03-p03-035-shelah-rank} below.
  \item \(\varphi(\bar x;\bar y)\) has the \emph{finite cover property}
        when for arbitrarily large \(n\) there are \(\bar a_0,\ldots,\bar
        a_{n-1}\) with \(\neg(\exists\bar x)\bigwedge_{k<n}\varphi(\bar
        x;\bar a_k)\) but \((\exists\bar x)\bigwedge_{k<n,k\neq l}\varphi(\bar
        x;\bar a_k)\) for every \(l<n\) (Definition~3.2, p.~312).
  \item \(\varphi(\bar x;\bar y)\) has the \emph{independence property}
        when for every \(n\) there are \(\bar a^0,\ldots,\bar a^{n-1}\) such
        that for every \(w\subseteq n\),
        \((\exists\bar x)\bigwedge_{k<n}\varphi(\bar x;\bar a^k)^{\text{if }(k\in w)}\)
        holds (Definition~4.1, pp.~316--317).
  \item \(\varphi(\bar x;\bar y)\) has the \emph{strict order property}
        when for every \(n\) there are \(\bar a^0,\ldots,\bar a^{n-1}\) such
        that for \(k,l<n\),
        \(\entails(\exists\bar x)[\neg\varphi(\bar x;\bar a^k)\wedge\varphi(\bar x;\bar a^l)]\)
        if and only if \(k<l\) (Definition~4.2, p.~317).
\end{enumerate}
\end{hbdefinition}

Shelah's independence property for \(\varphi(\bar x;\bar y)\) is
shattering of \(\{\bar a^0,\ldots,\bar a^{n-1}\}\) by the opposite formula,
so by \cref{thm:p03-w03-p03-035-ip-characterization}~(4) it is the IP of
\cref{def:p03-nip}.  His order property, realized for each \(r\) by some
\(\bar c_r\), says that \(\varphi(\bar c_r;\bar a^n)\) holds iff
\(r\leq n\); with \(\bar c_{i+1}\) in place of \(\bar c_i\) this is the order
property of \cref{def:p03-stability} for the same partitioned formula.  His
strict order property says that for \(l\leq k\) the set
\(\varphi(\mathfrak C;\bar a^l)\) is contained in
\(\varphi(\mathfrak C;\bar a^k)\) and for \(k<l\) it is not: the instances
form a strictly increasing chain of length \(n\) for every \(n\).  By
compactness this is the strict order property of
\cref{def:p03-w03-p03-035-sop}.

\begin{hbdefinition}{def:p03-w03-p03-035-shelah-rank}{Shelah's rank of \(\varphi\)-\(m\)-types}
Let \(\varphi(\bar x;\bar y)\) be a formula with \(\abs{\bar x}=m\).  For a
small set \(A\), a \emph{\(\varphi\)-\(m\)-type over \(A\)} is a set of
formulas \(\varphi(\bar x;\bar a)^{\varepsilon}\), with \(\bar a\) a tuple
from \(A\) of length \(\abs{\bar y}\) and \(\varepsilon\in\{0,1\}\), that is
realized in \(\mathfrak C\); it need not be complete, and \(S^m_\varphi(A)\)
is the set of complete ones.  By induction on the ordinal \(\alpha\) the sets
\(S^m_{\varphi,\alpha}(A)\) are defined for all small \(A\) at once
\cite[Definition~2.4, p.~299]{Shelah1971Stability}:
\begin{enumerate}
  \item \(S^m_{\varphi,0}(A)\) is the set of all \(\varphi\)-\(m\)-types over
        \(A\);
  \item \(S^m_{\varphi,\alpha+1}(A)\) is the set of those
        \(p\in S^m_{\varphi,\alpha}(A)\) such that for every finite
        \(q\subseteq p\) there are a small set \(B\supseteq A\) and a tuple
        \(\bar b\) from \(B\) with
        \(q\cup\{\varphi(\bar x;\bar b)\}\in S^m_{\varphi,\alpha}(B)\) and
        \(q\cup\{\neg\varphi(\bar x;\bar b)\}\in S^m_{\varphi,\alpha}(B)\);
  \item \(S^m_{\varphi,\delta}(A)=\bigcap_{\alpha<\delta}S^m_{\varphi,\alpha}(A)\)
        for a limit ordinal \(\delta\).
\end{enumerate}
For a \(\varphi\)-\(m\)-type \(p\) over \(A\), \(\operatorname{Rank}^m_\varphi p\)
is the greatest \(\alpha\) with \(p\in S^m_{\varphi,\alpha}(A)\), and is
\(\infty\) when \(p\in S^m_{\varphi,\alpha}(A)\) for every \(\alpha\); every
ordinal is counted below \(\infty\), and a set of such formulas that is not
realized in \(\mathfrak C\) receives rank \(-1\)
\cite[Definition~2.5, p.~300]{Shelah1971Stability}.  When \(\varphi\) and
\(m\) are fixed, \(\operatorname{Rank}p\) is written.
\index{rank!Shelah's local rank}
\end{hbdefinition}

Shelah works in a monster model and calls a set of formulas consistent when
it is realized there \cite[Lemma~2.7, p.~301]{Shelah1971Stability}; here
\(\mathfrak C\) plays that role, and the sets \(A\), \(B\) of the definition
are taken small so that ``realized'' and ``finitely realized'' agree by
saturation (\cref{thm:p03-w03-p03-031-saturation}).  Clause~(2) says that
every finite part of \(p\) can be split by one instance of \(\varphi\) into
two parts that both keep rank \(\alpha\).  That the greatest \(\alpha\)
exists and does not depend on the set \(A\) is part of the next lemma.

\begin{hblemma}{thm:p03-w03-p03-035-shelah-rank-basics}{proved=annex:p03-w03-p03-035-shelah-rank-basics}{Shelah's Lemmas~2.4 and 2.5: the rank is well defined and monotone}
Let \(\varphi(\bar x;\bar y)\) be a formula with \(\abs{\bar x}=m\), and let
\(p,q,q_1,q_2\) be sets of formulas \(\varphi(\bar x;\bar a)^{\varepsilon}\)
with parameters from small sets.
\begin{enumerate}
  \item \emph{(Lemma~2.4.)}  If \(p\) is a \(\varphi\)-\(m\)-type over \(A\)
        and \(A\subseteq B\) with \(B\) small, then for every \(\alpha\),
        \(p\in S^m_{\varphi,\alpha}(A)\) if and only if
        \(p\in S^m_{\varphi,\alpha}(B)\).  If \(\beta\leq\alpha\), then
        \(S^m_{\varphi,\alpha}(A)\subseteq S^m_{\varphi,\beta}(A)\).
        Hence \(\operatorname{Rank}p\) exists and depends only on \(p\).
  \item \emph{(Lemma~2.5\,(A).)}  If \(p\subseteq q\), then
        \(\operatorname{Rank}p\geq\operatorname{Rank}q\).
  \item \emph{(Lemma~2.5\,(B).)}  Every \(\varphi\)-\(m\)-type \(p\) has a
        finite subset \(q\) with \(\operatorname{Rank}q=\operatorname{Rank}p\).
  \item \emph{(Lemma~2.5\,(C).)}  If \(\operatorname{Rank}p=\alpha<\infty\),
        there are no \(q_1,q_2\supseteq p\) and \(\bar a\) with
        \(\operatorname{Rank}q_1=\operatorname{Rank}q_2=\alpha\),
        \(\varphi(\bar x;\bar a)\in q_1\) and
        \(\neg\varphi(\bar x;\bar a)\in q_2\).
  \item \emph{(Lemma~2.5\,(D).)}  A \(\varphi\)-\(m\)-type over \(A\) of
        rank \(\alpha<\infty\) has at most one extension in
        \(S^m_\varphi(A)\) of rank \(\alpha\).
  \item \emph{(Splitting.)}  If both \(p\cup\{\varphi(\bar x;\bar b)\}\) and
        \(p\cup\{\neg\varphi(\bar x;\bar b)\}\) have rank at least
        \(\alpha\), then \(\operatorname{Rank}p\geq\alpha+1\).
\end{enumerate}
\end{hblemma}

Shelah's Lemma~2.4 on p.~300 states the passage from \(A\) to \(B\); the
converse passage and the monotonicity in \(\alpha\) are the facts behind his
remark, after Definition~2.5, that the rank depends only on \(p\).  Item~(6)
is the form in which Lemma~2.5\,(C) is applied in the proof of his
Lemma~2.7 \cite[p.~302]{Shelah1971Stability}.  Items~(3) and (5) make a type
of rank below \(\infty\) determined, among extensions retaining its rank,
by a finite part of itself.  This rank-sensitive uniqueness is the counting
behind stability.

\begin{hblemma}{thm:p03-w03-p03-035-shelah-rank-trees}{proved=annex:p03-w03-p03-035-shelah-rank-trees}{Shelah's Lemmas~2.6--2.8: the rank and the binary tree}
Let \(\varphi(\bar x;\bar y)\) be a formula with \(\abs{\bar x}=m\).  For a
\(\varphi\)-\(m\)-type \(p\) and \(n<\omega\), let \(\Gamma_{p,n}\) be the set
of formulas \(\varphi(\bar x_\eta;\bar a)^{i}\) with
\(\varphi(\bar x;\bar a)^{i}\in p\) and \(\eta\in{}^{n}2\), together with
the formulas \(\varphi(\bar x_\eta;\bar y_{\eta\restriction k})^{\eta(k)}\)
with \(\eta\in{}^{n}2\) and \(k<n\).
\begin{enumerate}
  \item \emph{(Lemma~2.6.)}  If \(T\cup\Gamma(\varphi,m,\omega)\) is
        consistent, some \(\varphi\)-\(m\)-type has rank \(\infty\).
  \item \emph{(Lemma~2.7.)}  \(\operatorname{Rank}p\geq n\) if and only if
        \(\Gamma_{p,n}\) is realized in \(\mathfrak C\).
  \item \emph{(Lemma~2.8.)}  If there is no finite upper bound on the ranks
        of \(\varphi\)-\(m\)-types, then \(T\cup\Gamma(\varphi,m,\omega)\)
        is consistent.  If \(\operatorname{Rank}p\geq\omega\), then
        \(\operatorname{Rank}p=\infty\).
  \item The rank takes no value \(\alpha\) with
        \(\omega\leq\alpha<\infty\), and \(\operatorname{Rank}p\) equals the
        tree rank \(\mathrm R_\varphi(p)\) of
        \cref{def:p03-w03-p03-032-tree-rank}, read with the tuple \(\bar x\)
        as its object variables.
\end{enumerate}
\end{hblemma}

The set \(\Gamma_{p,n}\) is Shelah's own \cite[Lemma~2.7,
p.~301]{Shelah1971Stability}: a binary tree of height \(n\) of instances of
\(\varphi\), each of whose \(2^n\) leaves also satisfies \(p\).  For the
empty type it is \(\Gamma(\varphi,m,n)\) without \(T\).  Item~(4) compares
the rank with the tree rank of the stability-spectrum chapter, where the
rank is defined through these trees from the start; the ordinal induction of
\cref{def:p03-w03-p03-035-shelah-rank} gives the same values, and it never
stops strictly between \(\omega\) and \(\infty\).

\begin{hblemma}{thm:p03-w03-p03-035-shelah-tree}{proved=annex:p03-w03-p03-035-shelah-tree}{Shelah's Lemmas~2.2 and 2.3: from order to tree to many types}
Let \(\varphi(\bar x;\bar y)\) be a formula with \(\abs{\bar x}=m\).
\begin{enumerate}
  \item If \(\varphi\) has the order property, then
        \(T\cup\Gamma(\varphi,m,\omega)\) is consistent.
  \item If \(T\cup\Gamma(\varphi,m,\omega)\) is consistent and \(\lambda\) is
        an infinite cardinal, then some model of \(T\) contains a set \(A\)
        with \(\abs A\leq\lambda\) and \(\abs{S^m_\varphi(A)}>\lambda\); so
        \(\varphi\) is unstable in every infinite cardinal.
\end{enumerate}
\end{hblemma}

In (2), \(S^m_\varphi(A)\) is the set of complete \(\varphi\)-\(m\)-types over
\(A\) consistent with the elementary diagram of the model, as in Shelah's
Definition~2.1 on p.~297.  The same reading is used in~(A) and~(B) below.

\begin{hbtheorem}{thm:p03-w03-p03-035-shelah-29}{proved=annex:p03-w03-p03-035-shelah-29}{Shelah's Theorem~2.9: properties equivalent to instability of a formula}
For a formula \(\varphi(\bar x;\bar y)\) with \(\abs{\bar x}=m\) the
following are equivalent:
\textup{(A)} \(\varphi\) is unstable in every infinite cardinal;
\textup{(B)} \(\varphi\) is unstable in at least one infinite cardinal;
\textup{(C)} \(\varphi\) has the order property;
\textup{(D)} \(T\cup\Gamma(\varphi,m,\omega)\) is consistent;
\textup{(E)} not every \(\varphi\)-\(m\)-type has rank \(<\infty\);
\textup{(F)} there is no finite upper bound on the ranks of
\(\varphi\)-\(m\)-types.
Here \(\varphi\) is \emph{unstable in \(\lambda\)} when some model of \(T\)
contains a set \(A\) with \(\abs A\leq\lambda\) and
\(\abs{S^m_\varphi(A)}>\lambda\), and the rank is that of
\cref{def:p03-w03-p03-035-shelah-rank}.
\end{hbtheorem}

This is \cite[Theorem~2.9, p.~303]{Shelah1971Stability}.  Shelah proves it
on pp.~298--303 by two cycles of lemmas:
(A)\(\Rightarrow\)(B)\(\Rightarrow\)(C)\(\Rightarrow\)(D)\(\Rightarrow\)(A)
by Lemmas~2.1--2.3, and (D)\(\Rightarrow\)(E)\(\Rightarrow\)(F)\(\Rightarrow\)(D)
by Lemmas~2.4--2.8 on the rank.  The annex follows both cycles.  The step
from (B) to (C) is Shelah's Lemma~2.1 on p.~298, which rests on a theorem of
Erd\H{o}s and Makkai on a family of subsets with more members than the base
set; both are proved in the stability-spectrum chapter
(\cref{thm:p03-w03-p03-032-erdos-makkai,thm:p03-w03-p03-032-many-to-order}),
and the annex here transfers the cardinal reading of~(B) to the setting of
that proof.  The steps (C)\(\Rightarrow\)(D)\(\Rightarrow\)(A) are
\cref{thm:p03-w03-p03-035-shelah-tree}, (D)\(\Rightarrow\)(E) and
(F)\(\Rightarrow\)(D) are \cref{thm:p03-w03-p03-035-shelah-rank-trees}, and
(E)\(\Rightarrow\)(F) holds because rank \(\infty\) is at least every
\(n\).

Theorem~2.13 on p.~304 lifts this to
theories: \(T\) is unstable iff it is unstable in some \(\lambda=\lambda^{\abs
T}\), iff some formula is unstable, iff some formula \(\varphi(x;\bar y)\)
with \(x\) a single variable is unstable, iff some formula has an infinite
ordered sequence in the sense of \cref{def:p03-stability}.  Theorem~3.8 on
p.~312 shows that every unstable theory has the finite cover property.
Theorem~4.1 on p.~317 shows that \(T\) is unstable iff it has the
independence property or the strict order property;
\cref{thm:p03-w03-p03-035-ip-or-sop} below proves this theory-level
statement by Simon's route.  Theorems~4.2--4.4 on pp.~320--324 count
\(\varphi\)-types using \(\operatorname{Ded}(\lambda)\), the least cardinal
\(\mu\) for which no linear order of cardinality \(\mu\) has a dense subset
of cardinality \(\lambda\) (Definition~4.4, p.~320): if \(\varphi\) is
unstable and \(\lambda<\mu<\operatorname{Ded}(\lambda)\), some set of size at
most \(\lambda\) carries at least \(\mu\) complete \(\varphi\)-types
(Theorem~4.2), and a formula with
the independence property has \(2^\lambda\) types over some set of size
\(\lambda\) (Theorem~4.3).  Simon gives the corresponding counting
characterization of NIP, which depends on set-theoretic hypotheses
\cite[Proposition~2.69, p.~35]{Simon2015GuideManuscript}; the finite count of
\cref{thm:p03-w03-p03-035-sauer-shelah} needs none.

Two other Shelah papers of 1971 are distinct from this one.
\emph{Every two elementarily equivalent models have isomorphic ultrapowers}
\cite{Shelah1971Ultrapowers} constructs, without the generalized continuum
hypothesis, an ultrafilter over \(\lambda\) for which elementarily equivalent
models of power below the least \(\mu\) with \(2^\mu>\lambda\) have
isomorphic ultrapowers \cite[p.~224]{Shelah1971Ultrapowers}; it contains no
formula-property classification.  \emph{The number of non-isomorphic models
of an unstable first-order theory} \cite{Shelah1971NonIsomorphicModels} is
cited by the Stability paper, in its list of results, for the many-models
consequences of instability for pseudo-elementary classes
\cite[p.~287, items~G6 and G7]{Shelah1971Stability}; that spectrum statement
is \cref{thm:p03-unstable-spectrum} of the survey chapter.  The formula
properties above are the content of the Annals paper.

\section{Simple theories, dividing, and local character}
\label{sec:p03-w03-p03-035-simple}

Simplicity is \cref{def:p03-simple-theory}: local character of forking.
Dividing and forking are \cref{def:p03-dividing-forking}, read in
\(\mathfrak C\).  A complete type over \(B\) \emph{divides over}
\(A\subseteq B\) when it contains a formula that divides over \(A\).  The
basic calculus of these relations in every complete theory is
\cref{thm:p03-w03-p03-031-basic-calculus}; the facts below are the ones the
implications of \cref{sec:p03-w03-p03-035-intersections} use.

\begin{hblemma}{thm:p03-w03-p03-035-dividing-basics}{proved=annex:p03-w03-p03-035-dividing-basics}{Dividing: finite inconsistency and smaller bases}
Let \(A\subseteq C\) be small, \(\varphi(x;b)\) a formula.
\begin{enumerate}
  \item If \(\varphi(x;b)\) divides over \(A\), there are \(k<\omega\) and an
        \(A\)-indiscernible \((b_i:i<\omega)\) with \(b_0=b\) such that the
        indexed family \((\varphi(x;b_i):i<\omega)\) is \(k\)-inconsistent.
  \item If \(\varphi(x;b)\) divides over \(C\), it divides over \(A\).
  \item If \(\varphi(x;b)\) divides over \(A\), it forks over \(A\); a
        complete type over \(B\supseteq A\) that divides over \(A\) forks over
        \(A\).
\end{enumerate}
\end{hblemma}

\begin{hbdefinition}{def:p03-w03-p03-035-sop}{Strict order property}
A partitioned formula \(\psi(x;y)\) has the \emph{strict order property}
(SOP) in \(T\) when there are \(\abs y\)-tuples \((b_i:i<\omega)\) with
\(\psi(\mathfrak C;b_i)\subsetneq\psi(\mathfrak C;b_{i+1})\) for every
\(i<\omega\) \cite[Definition~2.65, p.~33]{Simon2015GuideManuscript}.  The
theory has SOP when some formula has it.
\index{strict order property}
\end{hbdefinition}

The strict order property asks for an infinite strictly increasing chain of
instances, which is stronger than an order property; it is the second
alternative in \cref{thm:p03-w03-p03-035-ip-or-sop}.

\begin{hbtheorem}{thm:p03-w03-p03-035-sop-not-simple}{proved=annex:p03-w03-p03-035-sop-not-simple}{The strict order property excludes simplicity}
If some formula has the strict order property in \(T\), then \(T\) is not
simple.  More precisely, with \(\kappa=\abs T^+\) there are a formula
\(\chi(u;v)\), parameters \((c_\alpha:\alpha<\kappa)\) and a complete type
\(p\) over \(B=\bigcup_{\alpha<\kappa}c_\alpha\) containing every
\(\chi(u;c_\alpha)\), such that \(\chi(u;c_\alpha)\) divides over
\(\bigcup_{\beta<\alpha}c_\beta\) for each \(\alpha\); hence \(p\) divides,
and forks, over every \(A\subseteq B\) with \(\abs A\leq\abs T\).
\end{hbtheorem}

The formula \(\chi\) says that \(u\) lies strictly between two parameters in
the preorder defined from the chain.  Dense linear orders are the model case
(\cref{thm:p03-w03-p03-035-dlo}).

\section{Stable, simple, NIP, and their proved intersections}
\label{sec:p03-w03-p03-035-intersections}

All statements of this section concern complete theories.

\begin{hbtheorem}{thm:p03-w03-p03-035-stable-nip}{proved=annex:p03-w03-p03-035-stable-nip}{Stable theories are NIP}
If a partitioned formula \(\varphi(x;y)\) has IP, then \(\varphi\) has the
order property of \cref{def:p03-stability}.  Hence every stable theory is NIP.
\end{hbtheorem}

\begin{hbtheorem}{thm:p03-w03-p03-035-ip-or-sop}{proved=annex:p03-w03-p03-035-ip-or-sop}{Instability splits into independence and strict order}
If \(T\) is unstable, then some formula has IP or some formula has SOP.
More precisely, if \(\varphi(x;y)\) has the order property and is NIP, then
for some \(n\) and some formula \(\theta(y;z)\) with \(\abs z=(n-2)\abs x\),
the formula \(\theta(y;z)\wedge\varphi(x;y)\), partitioned with object
variables \(y\) and parameter variables \(x z\), has SOP.
\end{hbtheorem}

This is \cite[Theorem~2.67, pp.~33--34]{Simon2015GuideManuscript}; Shelah's
Theorem~4.1(A) on p.~317 of \cite{Shelah1971Stability} states the converse as
well.  It follows from \cref{thm:p03-w03-p03-035-stable-nip} for IP, and for
SOP from the chain itself: if
\(\psi(\mathfrak C;b_i)\subsetneq\psi(\mathfrak C;b_{i+1})\) and
\(a_i\in\psi(\mathfrak C;b_{i+1})\setminus\psi(\mathfrak C;b_i)\), then
\(\psi(a_i;b_j)\) holds iff \(i<j\).

\begin{hbcorollary}{thm:p03-w03-p03-035-simple-nip-stable}{proved=annex:p03-w03-p03-035-simple-nip-stable}{Simple and NIP theories are stable}
A complete theory that is simple and NIP is stable.  Conversely a stable
theory is NIP (\cref{thm:p03-w03-p03-035-stable-nip}).
\end{hbcorollary}

Simon states the equivalence ``simple and NIP iff stable'' in his
introduction \cite[p.~9]{Simon2015GuideManuscript}.  The converse half also
needs \emph{stable implies simple}, which is
\cref{thm:p03-w03-p03-035-stable-simple} below.

\begin{hblemma}{thm:p03-w03-p03-035-tree-order}{proved=annex:p03-w03-p03-035-tree-order}{The tree property gives the order property}
Let \(\varphi(x;y)\) and \(k\geq2\).  Suppose there are
\((b_\eta:\eta\in\omega^{<\omega},\ \eta\neq\varnothing)\) such that for every
\(\nu\in\omega^\omega\) the set \(\{\varphi(x;b_{\nu\restriction l}):1\leq
l<\omega\}\) is realized, and for every \(\eta\in\omega^{<\omega}\) the set
\(\{\varphi(x;b_{\eta\frown\langle i\rangle}):i<\omega\}\) is
\(k\)-inconsistent.  Then \(\neg\varphi(x;y)\) has the order property; in
particular \(T\) is unstable.
\end{hblemma}

The hypothesis is Shelah's tree property for \(\varphi\) with \(k\).  The
argument chooses a branch through the tree one level at a time, avoiding
at each level the finitely many children satisfied by earlier witnesses.

\begin{hbtheorem}{thm:p03-w03-p03-035-stable-dividing}{proved=annex:p03-w03-p03-035-stable-dividing}{Stable theories have local character of dividing}
Let \(T\) be stable.  For every small set \(B\) and every complete type
\(p\) over \(B\) in finitely many variables there is \(A\subseteq B\) with
\(\abs A\leq\abs T\) such that \(p\) does not divide over \(A\).
\end{hbtheorem}

\Cref{thm:p03-w03-p03-035-stable-dividing} gives local character for
\emph{dividing}.  \Cref{def:p03-simple-theory} asks for local character of
\emph{forking}, which is the finite-disjunction closure of dividing.  The
route below does not pass through the theorem that forking and dividing
agree.  It uses a rank that drops along dividing formulas and does not rise
under finite disjunctions, so it drops along forking formulas as well; in a
stable theory the rank is finite, and choosing formulas of minimal rank
supplies the small base.

\begin{hbdefinition}{def:p03-w03-p03-035-drank}{Local dividing rank}
Let \(\delta(x;y)\) be a partitioned formula and \(k\geq1\).  For formulas
\(\theta(x)\) with parameters from \(\mathfrak C\), in the object variables
\(x\) of \(\delta\), define \(\operatorname{D}_{\delta,k}(\theta)\geq n\) by
induction on \(n<\omega\):
\begin{enumerate}
  \item \(\operatorname{D}_{\delta,k}(\theta)\geq0\) when \(\theta\) is
        realized in \(\mathfrak C\);
  \item \(\operatorname{D}_{\delta,k}(\theta)\geq n+1\) when there are
        \(\abs y\)-tuples \((d_l:l<\omega)\) such that the indexed family
        \((\delta(x;d_l):l<\omega)\) is \(k\)-inconsistent and
        \(\operatorname{D}_{\delta,k}(\theta\wedge\delta(x;d_l))\geq n\) for
        every \(l<\omega\).
\end{enumerate}
\(\operatorname{D}_{\delta,k}(\theta)\) is the largest \(n\) with
\(\operatorname{D}_{\delta,k}(\theta)\geq n\), it is \(-1\) when \(\theta\) is
not realized, and it is \(\infty\) when
\(\operatorname{D}_{\delta,k}(\theta)\geq n\) for every \(n\).
\index{dividing rank}
\end{hbdefinition}

The parameters \(d_l\) are not required to form an indiscernible sequence,
and nothing about the parameter set of \(\theta\) enters the definition;
this is what makes the rank monotone and blind to disjunctions.

\begin{hblemma}{thm:p03-w03-p03-035-drank-calculus}{proved=annex:p03-w03-p03-035-drank-calculus}{Calculus of the local dividing rank}
Fix \(\delta(x;y)\) and \(k\geq1\) and write \(\operatorname{D}\) for
\(\operatorname{D}_{\delta,k}\).  Let \(\theta,\theta',\theta_1,\ldots,
\theta_m\) be formulas in \(x\) with parameters and \(n<\omega\).
\begin{enumerate}
  \item \emph{(Downward closure.)}  If \(\operatorname D(\theta)\geq n+1\),
        then \(\operatorname D(\theta)\geq n\).
  \item \emph{(Monotonicity.)}  If
        \(\mathfrak C\entails\forall x(\theta\rightarrow\theta')\) and
        \(\operatorname D(\theta)\geq n\), then
        \(\operatorname D(\theta')\geq n\).
  \item \emph{(Invariance.)}  If \(\theta(x;z)\) is a formula without
        parameters and \(e\equiv e'\), then
        \(\operatorname D(\theta(x;e))\geq n\) if and only if
        \(\operatorname D(\theta(x;e'))\geq n\).
  \item \emph{(Disjunctions.)}  If
        \(\operatorname D(\theta_1\vee\cdots\vee\theta_m)\geq n\), then
        \(\operatorname D(\theta_j)\geq n\) for some \(j\).
  \item \emph{(Finiteness.)}  If \(T\) is stable, then
        \(\operatorname D(x=x)<\infty\), and hence
        \(\operatorname D(\theta)<\infty\) for every \(\theta\).
\end{enumerate}
\end{hblemma}

\begin{hbtheorem}{thm:p03-w03-p03-035-stable-simple}{proved=annex:p03-w03-p03-035-stable-simple}{Stable theories are simple}
Let \(T\) be stable.  For every small \(B\) and every complete type \(p\)
over \(B\) in finitely many variables, there is \(A\subseteq B\) with
\(\abs A\leq\abs T\) such that \(p\) does not fork over \(A\).
Consequently \(T\) is simple in the sense of \cref{def:p03-simple-theory},
for types in any number of variables and over parameter sets of any model
of \(T\).
\end{hbtheorem}

The proof takes, for each of the \(\abs T\) pairs \((\delta,k)\), a formula
of \(p\) of least \(\operatorname{D}_{\delta,k}\)-rank and lets \(A\) be the
set of their parameters.  A forking formula of \(p\) would imply a
disjunction of formulas dividing over \(A\); after coding the finitely many
disjuncts into one formula \(\delta\), each disjunct lowers the
\(\operatorname{D}_{\delta,k}\)-rank of the chosen formula, and so does their
disjunction, which contradicts minimality.

\begin{hbtheorem}{thm:p03-w03-p03-035-ominimal-nip}{proved=annex:p03-w03-p03-035-ominimal-nip}{O-minimal theories are NIP}
Let \(\Lang\) contain a binary relation symbol \(<\) and let \(T\) be a
complete \(\Lang\)-theory all of whose models are o-minimal in the sense of
\cref{def:p03-o-minimal}.  Then for every partitioned formula
\(\varphi(x;y)\) with \(\abs x=1\) there is \(m<\omega\) such that, in every
model and for every \(b\), no \(c_0<\cdots<c_m\) alternate in
\(\varphi(x;b)\); every such \(\varphi\) has VC-dimension at most \(m+1\);
and \(T\) is NIP.
\end{hbtheorem}

The hypothesis is on the theory.  For a single structure the passage to its
theory is a separate theorem.

\begin{hbtheorem}{thm:p03-w03-p03-035-kps}{cited=Simon2015GuideManuscript}{Transfer of o-minimality from a structure to its theory}
Let \(\mathcal M\) be a structure in a language containing \(<\), on which
\(<\) is a linear order, such that every subset of \(\dom\mathcal M\)
definable with parameters is a finite union of points and intervals.  Then
every model of \(\operatorname{Th}(\mathcal M)\) has the same property.
\end{hbtheorem}

Simon states this without proof at
\cite[Section~A.1.3, p.~148]{Simon2015GuideManuscript}; it is the theorem of
Knight, Pillay and Steinhorn on o-minimal theories.  It is not proved in this
chapter, and it is the only route used here from one o-minimal structure to
theory-level NIP.

\begin{hbcorollary}{thm:p03-w03-p03-035-structure-nip}{proved=annex:p03-w03-p03-035-structure-nip}{An o-minimal structure has an NIP theory, given the transfer}
If \(\mathcal M\) is o-minimal in the sense of \cref{def:p03-o-minimal}, then
\(\operatorname{Th}(\mathcal M)\) is NIP.  The proof uses the cited transfer
\cref{thm:p03-w03-p03-035-kps}.
\end{hbcorollary}

Simon's own sketch of ``o-minimal implies NIP'' uses o-minimality of the
theory, compactness, \cref{thm:p03-w03-p03-035-one-variable} and the
formula \(x\leq y\) \cite[Example~2.12, p.~17]{Simon2015GuideManuscript}; the
annex proof bounds alternation directly.

\begin{hbtheorem}{thm:p03-w03-p03-035-dlo}{proved=annex:p03-w03-p03-035-dlo}{Dense linear orders: o-minimal, NIP, unstable, not simple}
Let \(T_{\rm dlo}\) be the theory of \cref{def:p03-rn-dlo-data}, in the
language \(\{<\}\) with equality and nothing else.  Every model of
\(T_{\rm dlo}\) is o-minimal; \(T_{\rm dlo}\) is NIP; the formula \(x<y\) has
the order property, so \(T_{\rm dlo}\) is unstable; \(x<y\) has the strict
order property, so \(T_{\rm dlo}\) is not simple.
\end{hbtheorem}

\begin{hbcorollary}{thm:p03-w03-p03-035-rcf}{proved=annex:p03-w03-p03-035-rcf}{Real closed fields: o-minimal, NIP, unstable, not simple}
In the ordered-ring language of \cref{def:p03-w03-p03-037-rcf},
\(\mathsf{RCF}\) is NIP and not simple; it is o-minimal and unstable by
\cref{thm:p03-w03-p03-037-o-minimal,thm:p03-w03-p03-037-unstable}.
\end{hbcorollary}

The language matters.  The order of a real closed field is definable from
the ring operations (positive elements are nonzero squares), so the
ring-language theory is NIP and unstable as well; but quantifier elimination
needs \(<\) (\cref{thm:p03-w03-p03-037-ring-language}), and o-minimality as
defined requires the order symbol.  Expansions need not keep these
properties: the expansion of the real field by a predicate for the integers
is neither o-minimal nor NIP (\cref{thm:p03-w03-p03-037-expansion}).

O-minimality is strictly stronger than NIP for ordered structures: an
expansion of a dense linear order can keep NIP and lose o-minimality.  The
witness below adds one unary predicate.

\begin{hbdefinition}{def:p03-w03-p03-035-dense-predicate}{Dense linear order with a dense and codense predicate}
Let \(\Lang_P=\{<,P\}\) with \(<\) binary and \(P\) unary, and let
\(T_{\rm dp}\) consist of the axioms of a linear order on \(<\), the axioms
\(\forall a\exists x\,x<a\) and \(\forall a\exists x\,a<x\), and
\[
  \forall a\forall b\bigl(a<b\rightarrow\exists x(a<x\wedge x<b\wedge
  P(x))\bigr),\qquad
  \forall a\forall b\bigl(a<b\rightarrow\exists x(a<x\wedge x<b\wedge
  \neg P(x))\bigr).
\]
The structure \((\mathbb R,<,\mathbb Q)\), with \(P\) interpreted by the
rationals, is a model: between two reals there are a rational and an
irrational.
\end{hbdefinition}

\begin{hbtheorem}{thm:p03-w03-p03-035-dense-predicate}{proved=annex:p03-w03-p03-035-dense-predicate}{An NIP ordered theory that is not o-minimal}
\begin{enumerate}
  \item \(T_{\rm dp}\) has quantifier elimination, with \(\top\) and
        \(\bot\) counted as quantifier-free, and is complete.
  \item \(T_{\rm dp}\) is NIP.
  \item No model of \(T_{\rm dp}\) is o-minimal in the sense of
        \cref{def:p03-o-minimal}: the set defined by \(P(x)\) is not a finite
        union of points and intervals.
  \item The formula \(x<y\) has the order property and the strict order
        property, so \(T_{\rm dp}\) is unstable and not simple.
\end{enumerate}
\end{hbtheorem}

The reduct of a model of \(T_{\rm dp}\) to \(\{<\}\) is a model of
\(T_{\rm dlo}\), which is o-minimal (\cref{thm:p03-w03-p03-035-dlo}).  So the
ordered examples of this section mark the boundary cases separately: a
real closed field expands its dense order by ring operations and stays
o-minimal; expanding \((\mathbb R,<)\) by a dense and codense unary
predicate keeps NIP and loses o-minimality; expanding the real field by a
predicate for the integers loses both.  No general preservation of NIP or o-minimality under
expansions is asserted.

\begingroup\small
\setlength{\tabcolsep}{4pt}
\begin{longtable}{@{}p{0.20\textwidth}p{0.36\textwidth}p{0.36\textwidth}@{}}
\caption{Implications among dividing lines, with their proofs.}
\label{tab:p03-w03-p03-035-implications}\\
\toprule
Implication & Level and hypotheses & Proof or counterexample \\
\midrule
\endfirsthead
\caption[]{Implications among dividing lines (continued).}\\
\toprule
Implication & Level and hypotheses & Proof or counterexample \\
\midrule
\endhead
\midrule\multicolumn{3}{r}{\emph{continued on next page}}\\
\endfoot
\endlastfoot
IP \(\Rightarrow\) order property & one formula, any complete \(T\) &
\cref{thm:p03-w03-p03-035-stable-nip} \\
\addlinespace
stable \(\Rightarrow\) NIP & complete theory &
\cref{thm:p03-w03-p03-035-stable-nip} \\
\addlinespace
NIP \(\not\Rightarrow\) stable & complete theory &
\(T_{\rm dlo}\), \cref{thm:p03-w03-p03-035-dlo} \\
\addlinespace
unstable \(\Rightarrow\) IP or SOP & complete theory &
\cref{thm:p03-w03-p03-035-ip-or-sop} \\
\addlinespace
SOP \(\Rightarrow\) not simple & complete theory &
\cref{thm:p03-w03-p03-035-sop-not-simple} \\
\addlinespace
simple and NIP \(\Rightarrow\) stable & complete theory &
\cref{thm:p03-w03-p03-035-simple-nip-stable} \\
\addlinespace
stable \(\Rightarrow\) local character of dividing & complete theory, small
bases & \cref{thm:p03-w03-p03-035-stable-dividing} \\
\addlinespace
stable \(\Rightarrow\) simple & complete theory, local character of
forking, any model & \cref{thm:p03-w03-p03-035-stable-simple}, through the
rank of \cref{def:p03-w03-p03-035-drank} \\
\addlinespace
simple \(\not\Rightarrow\) stable & complete theory &
\(T_{\rm RG}\), \cref{thm:p03-w03-p03-035-random-graph} \\
\addlinespace
simple \(\not\Rightarrow\) NIP & complete theory &
\(T_{\rm RG}\), \cref{thm:p03-w03-p03-035-random-graph} \\
\addlinespace
o-minimal \(\Rightarrow\) NIP & theory all of whose models are o-minimal &
\cref{thm:p03-w03-p03-035-ominimal-nip} \\
\addlinespace
o-minimal structure \(\Rightarrow\) NIP theory & one structure
expanding a dense order & \cref{thm:p03-w03-p03-035-structure-nip}, using
the cited transfer \cref{thm:p03-w03-p03-035-kps} \\
\addlinespace
o-minimal \(\not\Rightarrow\) stable & language \(\{<\}\) &
\(T_{\rm dlo}\), \cref{thm:p03-w03-p03-035-dlo} \\
\addlinespace
NIP \(\not\Rightarrow\) o-minimal & ordered theory, language \(\{<,P\}\) &
\(T_{\rm dp}\), \cref{thm:p03-w03-p03-035-dense-predicate} \\
\addlinespace
expansions need not keep NIP or o-minimality & real field with a predicate
for \(\mathbb Z\) & \cref{thm:p03-w03-p03-037-expansion} \\
\bottomrule
\end{longtable}
\endgroup

The table refines the implication column of
\cref{tab:p03-classification-overview}: each row names the level at which
the implication is proved and the result that proves it.

\section{Random graphs as simple unstable examples}
\label{sec:p03-w03-p03-035-random-graph}

Let \(T_{\rm RG}\) be the theory of the random graph in the language
\(\{E\}\) (\cref{thm:p03-w03-p03-015-graph-types}): irreflexivity, symmetry,
and the extension axioms.  It is complete and has quantifier elimination.

\begin{hbtheorem}{thm:p03-w03-p03-035-random-graph-ip}{proved=annex:p03-w03-p03-035-random-graph-ip}{The edge relation shatters every finite set}
In every model of \(T_{\rm RG}\), every finite set of vertices is shattered
by \(E(x;y)\).  Hence \(E(x;y)\) has IP and \(T_{\rm RG}\) is not NIP.
\end{hbtheorem}

\begin{hbtheorem}{thm:p03-w03-p03-035-random-graph}{proved=annex:p03-w03-p03-035-random-graph}{The random graph is simple, unstable and has IP}
\(T_{\rm RG}\) is simple, unstable, and has the independence property.
Every complete type \(\operatorname{tp}(a/B)\) does not fork over the
finite set of entries of \(a\) that lie in \(B\).
\end{hbtheorem}

The forking computation is \cref{thm:p03-w03-p03-031-random-graph}, and the
order property of \(E\) is shown in \cref{def:p03-w03-p03-031-random-graph-example};
the annex assembles them with \cref{thm:p03-w03-p03-035-random-graph-ip}.
The random graph separates the classes in both directions that matter
here: it is simple but not stable, so simplicity does not imply stability;
and it is simple with IP, so simplicity does not imply NIP.  By
\cref{thm:p03-w03-p03-035-ip-or-sop} and
\cref{thm:p03-w03-p03-035-sop-not-simple} its instability must come from the
independence property, as it does.

\begin{hbwarning}{def:p03-w03-p03-035-random-graph-limits}{What the random graph does not show}
The random graph has no strict order property, by
\cref{thm:p03-w03-p03-035-sop-not-simple}, so it does not separate NIP from
the absence of SOP.  It is also \(\aleph_0\)-categorical, so its instability
is invisible in the countable spectrum; the unstable-spectrum theorem
\cref{thm:p03-unstable-spectrum} concerns cardinals above \(\abs T\) only.
\end{hbwarning}

\paragraph{Sources.}
Shelah's formula properties: \cite[\S\S 2--4]{Shelah1971Stability}.
NIP formulas and theories, examples, stability and the strict order property:
\cite[Sections~2.1, 2.3.2--2.3.4, pp.~13--36]{Simon2015GuideManuscript};
VC-dimension and the Sauer--Shelah lemma:
\cite[Section~6.1, pp.~85--87]{Simon2015GuideManuscript}; o-minimality:
\cite[Section~A.1.3, pp.~148--150]{Simon2015GuideManuscript}.
